\chapter{State Stratification and Lifetimes}
\label{ch:strata}

\chapterorient{Every piece of data a solver touches has a \emph{lifetime}---a
window of time over which it is alive, during which it may be read, perhaps
written, and after which it is gone. The configuration is set once and read
forever; the iterate evolves step by step; the scratch basis is born at a
restart and thrown away at the next; a single recurrence scalar is carried
across an inner loop and then forgotten. This chapter sorts \emph{all} of a
solver's state by lifetime. The sorting bins are the \emph{strata}, and getting
them right is what turns a tangle of instance variables into a design you can
reason about---and what turns the variant-handling of \cref{ch:operators} from a
discipline you must remember into a typing rule the calculus checks for you.}

\section{Why a solver's data has lifetimes}

In \cref{ch:records} we learned to bundle a solver's state into records and to
``update'' them without mutation, by building fresh copies. In \cref{ch:monad}
we learned to thread one such record---the evolving \code{SimState}---through a
chain of steps so the spreads stay out of the algorithm's way. Both chapters
treated ``state'' as one undifferentiated thing: a record that flows. This
chapter draws the lines \emph{inside} that flow.

The motivation is concrete. Open a real iterative solver and you find a dozen
named values living side by side: the operator it inverts, the tolerances it
chases, the current best guess, a counter, a convergence flag, a basis of search
directions, a small dense matrix of inner products, a couple of rotation
registers, one or two stray scalars. In the original C++ they are all
\emph{instance fields} of one solver object---all stored the same way, all
reachable from any method, all looking, syntactically, alike. But they are not
alike. They differ in the one property that matters most for understanding the
design: \emph{how long each one lives}.

\begin{itemize}[nosep]
  \item The operator and the tolerances are fixed when the solver is built and
  never change while it runs. They live for the \emph{whole solve}, read-only.
  \item The current guess, the counter, and the flag \emph{evolve}: each step
  produces new values from the old. They too live for the whole solve, but they
  are read \emph{and} written.
  \item The basis of search directions is scratch space. In a restarted method
  it is allocated fresh at each restart and discarded at the next. Its lifetime
  is \emph{one restart cycle}---shorter than the solve.
  \item A recurrence scalar like the one that forms a search-direction
  coefficient is threaded from each inner step to the next, then reborn from
  scratch the next time the solver is called. Its lifetime is \emph{one call's
  worth of inner steps}---narrower still.
\end{itemize}

Four different lifetimes, four kinds of data, all stored identically in the
source. \Cref{fig:nested} draws them as nested time-windows. The thesis of this
chapter is that making those windows \emph{visible in the type}---giving each
lifetime its own record, its own stratum---is not pedantry. It is the structural
move that makes the rest of the framework tractable.

\begin{figure}[t]
\centering
\begin{tikzpicture}[font=\footnotesize, x=1mm, y=1mm]
  % time axis
  \draw[->,simGray] (0,-3) -- (118,-3) node[right,font=\scriptsize]{time};
  \node[font=\scriptsize,text=simGray] at (4,-6) {build};
  \node[font=\scriptsize,text=simGray] at (60,-6) {the solve runs};
  % OpParams: spans the whole solve (built at construction, read forever)
  \fill[simBgGray] (10,18) rectangle (114,26);
  \draw[simGray,thick] (10,18) rectangle (114,26);
  \node[anchor=west] at (12,22) {\code{OpParams} \;\scriptsize(build-time, read-only)};
  % SimState: spans the whole solve, evolving
  \fill[simBgBlue] (16,8) rectangle (114,16);
  \draw[simBlue,thick] (16,8) rectangle (114,16);
  \node[anchor=west] at (18,12) {\code{SimState} \;\scriptsize(run-time, evolving)};
  % Krylov: one restart cycle (two of them shown)
  \fill[simBgGreen] (22,-1) rectangle (62,6);
  \draw[simGreen,thick] (22,-1) rectangle (62,6);
  \node[anchor=west,font=\scriptsize] at (24,2.5) {\code{Krylov}};
  \fill[simBgGreen] (66,-1) rectangle (106,6);
  \draw[simGreen,thick] (66,-1) rectangle (106,6);
  \node[anchor=west,font=\scriptsize] at (68,2.5) {\code{Krylov}};
  \node[font=\scriptsize,text=simGreen] at (44,-3.6) {restart cycle};
  \node[font=\scriptsize,text=simGreen] at (86,-3.6) {restart cycle};
  % recurrence carry: one apply / inner-loop sweep (tiny bars)
  \foreach \x in {24,30,36,42,48,54}{
    \fill[simBgRust] (\x,-9) rectangle (\x+4,-5.5);
    \draw[simRust,semithick] (\x,-9) rectangle (\x+4,-5.5);}
  \node[anchor=west,font=\scriptsize,text=simRust] at (60,-7.25) {recurrence carry (one inner sweep each)};
\end{tikzpicture}
\caption[The four strata as nested lifetimes]{The four strata of a solver's
state, drawn as time-windows. \code{OpParams} (the build-time configuration) is
alive for the whole run and never changes. \code{SimState} (the evolving state)
is alive for the whole run and changes every step. The ephemeral \code{Krylov}
workspace lives for a single restart cycle and is reborn at each restart. The
recurrence carry is narrower still---it threads one inner sweep and is reborn
each call. The windows \emph{nest}: each shorter-lived stratum sits inside a
longer-lived one. Reading a value from a window that has closed is a bug, and the
whole point of typing the strata separately is to make that bug impossible to
write.}
\label{fig:nested}
\end{figure}

\begin{keyidea}
A value's \emph{lifetime} is part of its type. Two values that are both
``state'' but live for different windows---the configuration that is read for the
whole solve, the scratch basis that dies at the next restart---are different
\emph{kinds} of state and belong in different records. When the calculus puts
each lifetime in its own stratum, ``who owns what, for how long'' stops being a
fact you keep in your head and becomes a fact the type signature states out
loud. This is the single idea the chapter develops; everything else is its
consequences.
\end{keyidea}

\section{The four strata}

We now name the four strata, each with its lifetime, its mutability, and the
real record of the framework that lives there. \Cref{tab:strata} is the map; the
subsections walk the rows.

\begin{table}[t]
\centering
\caption[The four strata of solver state]{The four strata, sorted by lifetime.
The first three are the common case of any iterative solve; the fourth appears
only when an inner loop threads a recurrence scalar. ``Mutability'' is from the
\emph{caller's} view: \code{SimState} is evolved by the solve but is a fresh
value each step, never mutated in place (\cref{ch:records}).}
\label{tab:strata}
\small
\setlength{\tabcolsep}{5pt}
\renewcommand{\arraystretch}{1.25}
\begin{tabular}{@{}p{20mm}p{30mm}p{20mm}p{30mm}@{}}
\toprule
Stratum & Lifetime & Mutability & Record \& example field \\
\midrule
\textbf{Build-time config} & the whole solve & read-only & \code{OpParams}: \code{tol}, \code{flexible} \\
\textbf{Run-time state} & the whole solve, evolving & evolved (fresh each step) & \code{SimState}: \code{x}, \code{it}, \code{converged} \\
\textbf{Ephemeral workspace} & one restart cycle & internal scratch & \code{Krylov}: basis \code{V}, direction \code{p} \\
\textbf{Recurrence carry} & one inner sweep, reborn per call & threaded scalar & \code{PrevCarry}: \code{beta\_prev} \\
\bottomrule
\end{tabular}
\end{table}

\subsection{Stratum one: \texorpdfstring{\code{OpParams}}{OpParams}---build-time configuration}

The first stratum is everything the solver is \emph{configured} with: the
operator it applies, the preconditioner, the tolerances, the iteration limit,
and---crucially---the \emph{variant selectors} that pick which flavour of the
algorithm to run (left or right preconditioning, which Gram--Schmidt variant,
flexible or ordinary). All of it is captured once, at the moment the solver is
constructed, and none of it changes while the solve runs. Its lifetime is the
whole solve; its mutability, from anyone's view, is read-only.

The framework's record for this stratum is \code{OpParams}. Its shape, in the
record syntax of \cref{ch:records}:

\begin{lstlisting}[style=l4style]
OpParams = {
  -- constructed-operator surfaces (the kernel touches OpParams ONLY through these)
  T          : ConstructedOp,      -- the apply surface (preconditioned operator)
  eps        : Convergence,        -- the stopping-predicate surface

  -- variant selectors (closed over by the surfaces above; NOT read by the kernel body)
  pc_side    : PreconditionerSide, -- LEFT | RIGHT
  gs_orthog  : Orthogonalization,  -- MGS | CGS | CGS2
  flexible   : Bool,               -- FGMRES vs GMRES

  -- termination knobs (close into eps)
  max_dim    : Int,                -- restart subspace dimension
  max_it     : Int,
  rel_tol    : Scalar,
  abs_tol    : Scalar
}                                  -- the WHOLE record is readonly
\end{lstlisting}

Two things deserve immediate notice. First, the whole record is marked
\code{readonly}: there is no field a step writes. That is the defining property
of the stratum---build-time configuration is, by definition, the data that is
fixed before the loop and read during it. Second, the record is in two layers.
The \emph{variant selectors} (\code{pc\_side}, \code{gs\_orthog},
\code{flexible}) are raw choices; but the kernel never reads them directly.
Instead they are folded, at construction, into the \emph{constructed-operator
surfaces} (\code{T}, \code{eps})---closures that have already ``baked in'' the
variant choice. The per-step kernel touches \code{OpParams} only through those
surfaces. We will see in \cref{sec:visible} that this two-layer shape is exactly
what makes variant handling a checkable invariant rather than a discipline. The
fields trace to the immutable instance fields of Palace's \code{IterativeSolver}
hierarchy, which the framework un-mixes into this one read-only record.

\subsection{Stratum two: \texorpdfstring{\code{SimState}}{SimState}---run-time evolving state}

The second stratum is the state the solve \emph{evolves}: the externally-visible
quantities a caller observes after the solve returns. This is the record
\cref{ch:records} drew as cards and \cref{ch:monad} threaded through the
\code{Solve} monad. Its lifetime is also the whole solve, but unlike
\code{OpParams} it is read \emph{and} written---a new value every step. Its
shape is the familiar five fields:

\begin{lstlisting}[style=l4style]
SimState = {
  x           : Tensor[N],   -- the current iterate (the solve's primary product)
  it          : Int,         -- iteration count
  converged   : Bool,        -- convergence flag
  final_res   : Scalar,      -- final residual
  initial_res : Scalar       -- residual captured at solve entry
}
\end{lstlisting}

\code{SimState} is run-time-evolved \emph{but immutable}: ``evolves'' here means
exactly what it meant in \cref{ch:records}---each step builds a \emph{fresh}
\code{SimState}, like the previous one but for the fields that changed, and the
\code{Solve} monad of \cref{ch:monad} threads the fresh value forward. From the
\emph{caller's} viewpoint \code{SimState} is read-only (the solve hands back a
new value rather than reaching into theirs); from the \emph{solve's} viewpoint
it is the one bundle that ticks forward step by step. Notice how little actually
changes per step: the iteration counter \code{it} increments---typically the
kernel's \emph{sole} per-step monadic effect---while the heavy iterate \code{x}
is folded only at restart-cycle boundaries, and the exit fields
(\code{converged}, \code{final\_res}) are written only as the solve terminates.
This is the single-mutation-point observation we flagged in \cref{ch:records},
now made structural by the stratum boundary.

\Cref{fig:simstatecard} is the record card, with every field tagged
\emph{run-time} to mark the stratum. The decisive fact about \code{SimState}---
the one that makes it the monad's state type---is that its shape is
\emph{uniform across algorithms}. Conjugate gradients, GMRES, flexible GMRES,
and the Chebyshev smoother all evolve exactly these five fields. They differ
wildly in their workspace and their configuration, but the externally-visible
state they thread is one shape. That uniformity is what lets the \code{Solve}
monad's effect domain be exactly \code{SimState}: the monadic coordination of
\cref{ch:monad} does not vary by algorithm, because the thing it coordinates
does not. The fields mirror the \code{mutable} solve-statistics of Palace's base
class---\code{converged}, \code{initial\_res}, \code{final\_res}, the iteration
count---which the \code{mutable} keyword itself marks as the run-time stratum.

\begin{figure}[t]
\centering
\begin{tikzpicture}[font=\footnotesize]
  \node[draw=simBlue, fill=simBgBlue, rounded corners=3pt, thick,
        minimum width=86mm, minimum height=44mm, inner sep=6pt] (card) at (0,0) {};
  \node[anchor=north, font=\scriptsize\bfseries, text=simBlue]
    at ($(card.north)+(0,-1.5mm)$) {\code{SimState} \;\textnormal{\scriptsize(stratum 2: run-time)}};
  % rows
  \node[anchor=west] at ($(card.north west)+(4mm,-9mm)$)
    {\code{x}\,:\ \(\mathsf{Tensor}[N]\)\quad\scriptsize\itshape current iterate};
  \node[anchor=west] at ($(card.north west)+(4mm,-16mm)$)
    {\code{it}\,:\ \(\mathsf{Int}\)\quad\scriptsize\itshape iteration count};
  \node[anchor=west] at ($(card.north west)+(4mm,-23mm)$)
    {\code{converged}\,:\ \(\mathsf{Bool}\)\quad\scriptsize\itshape convergence flag};
  \node[anchor=west] at ($(card.north west)+(4mm,-30mm)$)
    {\code{final\_res}\,:\ \(\mathsf{Scalar}\)\quad\scriptsize\itshape final residual};
  \node[anchor=west] at ($(card.north west)+(4mm,-37mm)$)
    {\code{initial\_res}\,:\ \(\mathsf{Scalar}\)\quad\scriptsize\itshape entry residual};
  % run-time tags down the right edge
  \foreach \y in {9,16,23,30,37}{
    \node[anchor=east, font=\scriptsize, text=simRust]
      at ($(card.north east)+(-3mm,-\y mm)$) {\textsf{run-time}};}
  % rules
  \foreach \y in {12.5,19.5,26.5,33.5}{
    \draw[simGray!35] ($(card.north west)+(3mm,-\y mm)$) -- ($(card.north east)+(-3mm,-\y mm)$);}
\end{tikzpicture}
\caption[The \code{SimState} record card]{The \code{SimState} record, with all
five fields tagged \emph{run-time} to mark the stratum. This is the only state a
caller observes after the solve returns; it is the value the \code{Solve} monad
of \cref{ch:monad} threads. Its shape is identical for conjugate gradients,
GMRES, flexible GMRES, and Chebyshev---a single five-field card across every
algorithm---which is precisely what lets one monad coordinate them all.}
\label{fig:simstatecard}
\end{figure}

\subsection{Stratum three: ephemeral \texorpdfstring{\code{Krylov}}{Krylov} workspace}

The third stratum is \emph{scratch}: working data whose entire life is internal
to the solve and which is thrown away on return. For a GMRES-shaped method this
is the \code{Krylov} bundle---the orthonormal basis \code{V} of search
directions, the small dense Hessenberg matrix \code{H}, the least-squares
right-hand side \code{s}, the Givens rotation registers \code{cs}/\code{sn}. For
conjugate gradients it is the search direction \code{p} and the residual vector
\code{r}. Whatever the method, this stratum holds the algorithm's gears: large
enough to dominate memory, never visible to the caller.

\begin{lstlisting}[style=l4style]
-- GMRES / FGMRES workspace
Krylov = {
  V  : Basis,    -- orthonormal Krylov basis (the search directions)
  H  : Dense,    -- Hessenberg matrix
  s  : Vec,      -- least-squares right-hand side
  cs : Vec, sn : Vec,  -- Givens rotation registers
  j  : Int       -- position within the current restart cycle
}
\end{lstlisting}

Two properties set this stratum apart, and both are about lifetime. First, its
lifetime is \emph{shorter than the solve}. In a restarted method the basis
\code{V} fills up over an inner cycle and is then discarded: at each restart the
\code{Krylov} bundle is \emph{reborn}---a fresh bundle, not the old one with its
position counter reset. The window in \cref{fig:nested} that says ``one restart
cycle'' is this stratum's, and there are as many such windows as there are
restarts. Second---and this is the subtle structural point---\code{Krylov} is
\emph{not threaded by the monad}. It is threaded as a \emph{plain value},
passed from inner step to inner step by ordinary function application, never
entering the \code{Solve} monad's state.

Why keep it out of the monad? Because the monad threads \code{SimState}, and
\code{SimState} is exactly the state that survives the whole solve. The
\code{Krylov} bundle does \emph{not} survive the whole solve---it is reborn at
every restart. To thread it through the monad would be to claim it has the
lifetime of \code{SimState}, which is false. Its born-at-restart,
discarded-at-restart lifetime \emph{defeats} monadic encoding: the monad is the
wrong tool for a value whose window is a restart cycle, so we thread it as a
plain returned value instead. This is the first place where the lifetime of a
value dictates the \emph{mechanism} that carries it, and it is why the strata are
not a bookkeeping nicety but a design constraint.

\subsection{Stratum four: the recurrence carry}

The first three strata are the whole story for many solvers. But some carry a
\emph{fourth} kind of state, narrower than all three, and it is worth a stratum
of its own because mis-housing it is a real bug class.

Consider the conjugate-gradient direction update. Each step forms the new search
direction from the old one with a coefficient $\beta_k/\beta_{k-1}$, where
$\beta_k = (\mathsf{B}\vr_k, \vr_k)$ is this step's preconditioned residual inner
product and $\beta_{k-1}$ was last step's. To compute step $k$ you need a value
from step $k-1$: the recurrence \emph{threads}. The Chebyshev smoother has the
same shape, carrying a scalar $\rho_{k-1}$ through its inner polynomial
recurrence $\rho_k = 1/(2\theta/\delta - \rho_{k-1})$.

That single threaded scalar is the fourth stratum. The framework records it as
\code{PrevCarry}:

\begin{lstlisting}[style=l4style]
PrevCarry = { beta_prev: Scalar }   -- CG: the prior step's (Br, r)
-- (GMRES variant: PrevCarry = { H_prev: Scalar })
\end{lstlisting}

Its lifetime is the diagnostic. It is \emph{threaded across the inner loop}---
each step reads the previous step's value, so it has a genuine cross-iteration
data dependence---but it is \emph{reborn at each top-level call}: the next time
the solver runs, the recurrence restarts from its base case. ``One call's worth
of inner steps, then gone.'' That window is narrower than \code{OpParams} (which
lives across \emph{all} calls) and narrower than an ordinary scratch temporary
(which has no cross-iteration dependence at all). It is its own thing, and
\cref{fig:fourstrata} sorts a value into it by exactly these two questions.

\begin{figure}[t]
\centering
\resizebox{\textwidth}{!}{%
\begin{tikzpicture}[font=\footnotesize, node distance=7mm and 10mm]
  \tikzset{q/.style={draw=simBlue, fill=simBgBlue, rounded corners=2pt, thick,
      align=center, inner sep=4pt, minimum width=38mm, font=\footnotesize},
    leaf/.style={draw=simGray, fill=white, rounded corners=2pt, semithick,
      align=center, inner sep=4pt, minimum width=28mm, font=\scriptsize}}
  \node[q] (q1) {Captured at construction\\and read forever?};
  \node[leaf, below left=of q1] (op) {\textbf{OpParams}\\\scriptsize build-time config};
  \node[q, below right=of q1] (q2) {Read \emph{and written}, and\\survives the whole solve?};
  \node[leaf, below left=of q2] (sim) {\textbf{SimState}\\\scriptsize run-time state};
  \node[q, below right=of q2] (q3) {Has a cross-iteration\\recurrence dependence?};
  \node[leaf, below left=of q3] (kry) {\textbf{Krylov}\\\scriptsize ephemeral workspace};
  \node[leaf, below right=of q3] (car) {\textbf{PrevCarry}\\\scriptsize recurrence carry};
  \draw[flow] (q1) -- node[above left,font=\scriptsize]{yes} (op);
  \draw[flow] (q1) -- node[above right,font=\scriptsize]{no} (q2);
  \draw[flow] (q2) -- node[above left,font=\scriptsize]{yes} (sim);
  \draw[flow] (q2) -- node[above right,font=\scriptsize]{no} (q3);
  \draw[flow] (q3) -- node[above left,font=\scriptsize]{no} (kry);
  \draw[flow] (q3) -- node[above right,font=\scriptsize]{yes} (car);
\end{tikzpicture}}
\caption[Sorting a value into its stratum]{A decision tree for placing a piece of
solver state. The two questions that matter are about \emph{lifetime}: does the
value survive the whole solve (and is it written), and does it carry a
cross-iteration recurrence? Build-time configuration is captured once and read
forever (\code{OpParams}); evolving caller-visible state survives the solve
(\code{SimState}); scratch with no cross-step dependence is ephemeral workspace
(\code{Krylov}); a threaded recurrence scalar that is reborn each call is the
carry (\code{PrevCarry}).}
\label{fig:fourstrata}
\end{figure}

The fourth stratum is easy to confuse with two of the others, and the confusions
are both bugs. It is \emph{not} build-time configuration: if you stored
$\beta_{k-1}$ in \code{OpParams}, it would persist between calls and corrupt the
next solve, which must restart its recurrence from scratch. And it is \emph{not}
an ordinary ephemeral temporary: those (the residual vector recomputed each step,
the matrix-vector products) have no dependence on the previous iteration, whereas
the carry's whole nature is that step $k$ reads step $k-1$. We return to all
three confusions as pitfalls in \cref{sec:mistakes}.

\section{Why lifetimes must be visible}
\label{sec:visible}

We have four strata and four records. The question this section answers is
\emph{why bother}---why not store everything as instance fields the way the C++
does, and keep the lifetimes in our heads? The answer is that making the
lifetimes visible in the type buys two guarantees that are otherwise only
hopes.

\subsection{You cannot persist ephemeral data past its window}

The first guarantee is a safety property. When the ephemeral \code{Krylov}
workspace is its own value, threaded as a plain value and reborn at each restart,
there is \emph{no place} to accidentally read it after its window closes. The
type says: at the restart boundary, the old \code{Krylov} is gone and a fresh one
takes its place. A step that tried to read a basis vector from before the restart
would be reading from a value that no longer exists---and because the fresh
bundle is a genuinely new value (\cref{ch:records}), there is no aliased box
holding the stale data for it to find.

Contrast the mutable world. There, the basis \code{V} is an instance field that
outlives any one restart cycle; ``reborn at restart'' is implemented by
\emph{overwriting} it in place, and nothing stops a stray method from reading a
slot that belongs to the previous cycle. The lifetime is real but invisible---a
fact about when the code happens to overwrite \code{V}, not a fact the type
enforces. The stratum makes the invisible visible: a restart \emph{produces a
new bundle}, so reading across the boundary is reading a value that was never
built.

\subsection{Variant absorption becomes a typing invariant}

The second guarantee is the deeper one, and it is the bridge to
\cref{ch:operators}. Recall the two-layer shape of \code{OpParams}: variant
selectors on one layer, constructed-operator surfaces on another, and the rule
that \emph{the per-step kernel touches \code{OpParams} only through the
surfaces}, never through the raw selectors. Why is that rule worth stating?

Because it makes \emph{variant absorption} checkable. A single kernel is supposed
to run all the variants---left and right preconditioning, the three
Gram--Schmidt flavours, flexible and ordinary---without itself branching on which
variant it is. The variant choice is supposed to be ``absorbed'' at construction,
into the surfaces, so the kernel's control flow is uniform. In a procedural
solver that uniformity is a \emph{discipline}: the programmer must remember not
to write \code{if (pc\_side == LEFT)} in the inner loop. Nothing checks it.

With the strata, it is a \emph{typing invariant}. If \code{OpParams} is
\code{readonly}, and the variant selectors are reachable from the kernel only
\emph{through} the constructed surfaces (because the kernel is handed the
surfaces, not the raw record), then the kernel \emph{provably cannot} branch on a
selector---there is no expression it could write to read one. Variant absorption
stops being ``a thing we are careful to do'' and becomes ``a thing the type
forbids us from violating.'' That is the entire payoff of separating the
build-time stratum from the run-time one, and \cref{ch:operators} cashes it in to
show one kernel covering an entire family of algorithms.

\begin{keyidea}
Visible lifetimes turn two correctness properties from discipline into typing.
\emph{Safety:} because ephemeral workspace is a fresh value reborn at each
restart, there is no stale box to read across the boundary---persisting it past
its window is unwriteable, not merely discouraged. \emph{Variant absorption:}
because \code{OpParams} is read-only and the kernel reaches the variant selectors
only through the constructed surfaces, the kernel \emph{cannot} branch on a
variant---the absorption invariant of \cref{ch:operators} is checked by the
types, not remembered by the programmer. Both guarantees are consequences of one
move: giving each lifetime its own stratum.
\end{keyidea}

\section{Build-time versus run-time: the two great strata}

Step back from the four-way split and a coarser, even more important division
appears. Group the strata by \emph{when their work runs}, and there are exactly
two camps:

\begin{itemize}[nosep]
  \item \textbf{Build-time work} runs \emph{once}, before the loop. Constructing
  the operator surfaces, factoring a preconditioner, installing pointers, parsing
  the configuration record, querying the function-space hierarchy---all of it
  happens at solver construction and is done by the time the first iteration
  begins. This is the world of \code{OpParams}.
  \item \textbf{Run-time work} runs \emph{every step}, inside the loop. Applying
  the operator, applying the preconditioner, bumping the counter, computing a
  residual norm, folding the iterate. This is the world of \code{SimState}, the
  \code{Krylov} workspace, and the carry.
\end{itemize}

\Cref{fig:buildrun} draws the two as a timeline: a setup pass, then the loop.
This build-time/run-time line is the same one we drew for the matrix-free
operator in \cref{ch:matrixfree}, where the operator's \emph{setup}---assembling
the element-local quadrature data once---is split from its \emph{apply}---the
per-call tensor contraction that runs every iteration. There, as here, the setup
is build-time and the apply is run-time; the two were separated for exactly the
reason the strata separate them now.

\begin{figure}[t]
\centering
\begin{tikzpicture}[font=\footnotesize, x=1mm, y=1mm]
  % timeline
  \draw[->,simGray] (0,0) -- (122,0) node[right,font=\scriptsize]{time};
  % build-time pass
  \fill[simBgGray] (4,4) rectangle (40,16);
  \draw[simGray,thick] (4,4) rectangle (40,16);
  \node[align=center,font=\scriptsize] at (22,13) {\textbf{build-time pass} (once)};
  \node[align=center,font=\scriptsize,text=simGray] at (22,8.5)
    {construct surfaces\\factor preconditioner\\parse config};
  \node[font=\scriptsize,text=simGray,anchor=north] at (22,3) {\code{OpParams} born here};
  % run-time loop
  \fill[simBgBlue] (48,4) rectangle (118,16);
  \draw[simBlue,thick] (48,4) rectangle (118,16);
  \node[align=center,font=\scriptsize] at (83,13) {\textbf{run-time loop} (every step)};
  % loop iterations as little boxes
  \foreach \x in {52,62,72,82,92,102}{
    \draw[simBlue,fill=white,semithick] (\x,6) rectangle (\x+7,10.5);
    \node[font=\tiny] at (\x+3.5,8.25) {apply};}
  \node[font=\scriptsize,text=simBlue,anchor=north] at (83,3)
    {\code{SimState} / \code{Krylov} evolve here};
  % the boundary
  \draw[->,flow] (40,10) -- (48,10);
  \node[font=\scriptsize,text=simRust,align=center] at (44,16) {loop\\entry};
\end{tikzpicture}
\caption[The build-time / run-time timeline]{Solver work splits into a setup
pass that runs once and a loop that runs every step. Build-time work---building
the operator surfaces, factoring the preconditioner, parsing configuration---is
done before the first iteration and populates \code{OpParams}. Run-time
work---the per-step applies, counter bumps, residual computations---runs inside
the loop and evolves \code{SimState} and the ephemeral workspace. This is the
same split as the matrix-free operator's setup-versus-apply of
\cref{ch:matrixfree}.}
\label{fig:buildrun}
\end{figure}

The division matters for a reason beyond tidiness, and it connects to the very
purpose of the layered framework. \emph{Only run-time work lifts to the
tensor-field form of \cref{ch:tensors}.} When we re-express a solve as global
tensor-field operations---the whole point of the climb toward the abstract
calculus---it is the per-element, per-iteration work that has structure to lift:
an \code{apply} is a tensor contraction, a residual is a reduction, a counter
bump is a scalar increment. Build-time work has no per-element structure to lift;
it \emph{sets up} the loop rather than \emph{being} a loop. So the build-time
stratum is recorded as the scaffolding that materializes the operator graph, and
it is the run-time stratum---and only the run-time stratum---that participates in
the lift. The build-time/run-time line is thus not only a lifetime distinction;
it is the line between the work that becomes a tensor-field expression and the
work that prepares for it.

In the monad of \cref{ch:monad} this lands as a clean structural rule:
build-time work sits in the \code{do}-block's \emph{setup} prologue, before the
loop; run-time work sits in the \emph{body} that the loop iterates. Schematically,

\begin{lstlisting}[style=l4style]
solve = do
  setup <- build_operators(config)   -- BUILD-TIME: once, before the loop
  iterate_while(not_converged) (\s ->
     step(setup, s))                  -- RUN-TIME: every iteration
\end{lstlisting}

The \code{setup} line runs once and produces the \code{OpParams}-stratum
surfaces; the \code{step} inside \code{iterate\_while} runs every iteration and
evolves the run-time strata. Reading the \code{do}-block, you can see the
lifetime of every value by \emph{where} it appears: above the loop, it is
build-time; inside, it is run-time.

\begin{workedexample}{Sorting a conjugate-gradient solve, variable by variable}{cgsort}
We take a complete preconditioned conjugate-gradient solve and sort
\emph{every} variable it names into its stratum. The exercise is the whole
chapter in miniature: given a pile of state, place each value by its lifetime.

\textbf{The variables.} A PCG solve holds: the system operator $\mathsf{A}$ and
preconditioner $\mathsf{B}$; the relative and absolute tolerances and the
iteration limit; the current iterate $\vx$; the iteration count and the
convergence flag; the residual vector $\vr$ and the preconditioned residual
$\mathsf{B}\vr$; the search direction $\vp$; the step lengths $\alpha$ formed
each step; and the recurrence scalar $\beta = (\mathsf{B}\vr,\vr)$, whose
\emph{previous} value $\beta_{k-1}$ feeds the direction-update coefficient.

\tcblower
\textbf{The sort.} Ask of each value the two lifetime questions of
\cref{fig:fourstrata}.

\begin{itemize}[nosep]
  \item \(\mathsf{A}\), \(\mathsf{B}\), \code{rel\_tol}, \code{abs\_tol},
  \code{max\_it} --- captured at construction, read forever, never written.
  \textbf{Stratum 1, \code{OpParams}.} (The operators are reached through the
  constructed \code{T} surface; the tolerances close into \code{eps}.)
  \item \(\vx\), \code{it}, \code{converged}, the residual proxies --- evolved
  every step, survive the whole solve, visible to the caller.
  \textbf{Stratum 2, \code{SimState}.}
  \item \(\vr\), \(\mathsf{B}\vr\), \(\vp\), and the per-step \(\alpha\) ---
  scratch, recomputed each step, with no value carried from one step to the next;
  discarded on return. \textbf{Stratum 3, ephemeral \code{Krylov}.}
  \item \(\beta_{k-1}\) --- the one scalar with a genuine cross-step dependence
  (\(\vp_k = \vr_k + (\beta_k/\beta_{k-1})\,\vp_{k-1}\)), threaded inner step to
  inner step but restarting from its base case each new solve.
  \textbf{Stratum 4, the carry \code{PrevCarry}.}
\end{itemize}

\textbf{The trap, named.} Note that $\vp$ and $\beta_{k-1}$ both ``come from the
previous step,'' and the lazy sort would lump them together. They are in
different strata. The direction $\vp$ is rebuilt each step from current
quantities (it has a \emph{recurrence in form} but we treat the rebuilt vector as
ephemeral workspace); the \emph{scalar} $\beta_{k-1}$ is the value with the true
cross-iteration data dependence that the carry exists to thread. The
distinguishing question---``does step $k$ read a \emph{value} from step
$k-1$?''---is what separates stratum 4 from stratum 3, and it is exactly the
bottom question of \cref{fig:fourstrata}.
\end{workedexample}

\section{The book's real records, defined}

We have met all five records in passing. We now define each one as a card of
named fields, with the stratum each field lives in stated outright. This section
is the reference the rest of Part~III points back to.

\subsection{\texorpdfstring{\code{SolveResult}}{SolveResult}---what a step returns}

A single step does not return a bare \code{SimState}; it returns a bundle. The
framework calls the bundle \code{SolveResult}, written with the \code{Solve}
name:

\begin{lstlisting}[style=l4style]
SolveResult = Solve { sim: SimState, krylov: Krylov, outputs: StepOutputs }
-- (Form B, when the carry is threaded explicitly:)
--   Solve { sim: SimState, krylov: Krylov, outputs: StepOutputs, carry: PrevCarry }
\end{lstlisting}

Read it field by field, by stratum:

\begin{itemize}[nosep]
  \item \code{sim} --- the next \code{SimState}, stratum 2. This is the one field
  the step delivers \emph{through the monad}: it is the monadic-effect product,
  the state transition the \code{Solve} monad of \cref{ch:monad} threads.
  Typically the only thing that changed is the \code{it} counter.
  \item \code{krylov} --- the next ephemeral workspace, stratum 3, returned as a
  \emph{plain value} (because, as we saw, its restart-bounded lifetime defeats
  monadic threading).
  \item \code{outputs} --- the demand-prunable per-step readouts, a fresh bundle
  each step; defined just below.
  \item \code{carry} --- present only in Form B, the recurrence carry of stratum
  4 (\code{PrevCarry}), threaded back into the next step.
\end{itemize}

There is a deliberate overload here, and \cref{ch:monad} warned of it:
\code{Solve} names \emph{both} the monad (the effect that threads \code{SimState})
\emph{and} this returned record (the braces). They are different things sharing a
spelling: the monad answers ``how is state threaded?''; \code{SolveResult}
answers ``what record does a step yield?''. The \code{sim} field is where they
touch---it is the effect's state product---while the other fields are the plain
returned value. \Cref{ch:monad} owns the threading; this stratum chapter owns the
record's fields and their lifetimes.

\subsection{\texorpdfstring{\code{StepOutputs}}{StepOutputs}---demand-prunable per-step readouts}

The \code{outputs} field is its own record, and it is the textbook home of the
optional-field, demand-pruning idea of \cref{ch:records}:

\begin{lstlisting}[style=l4style]
StepOutputs = {
  residual_norm:    Scalar,        -- the step's residual proxy
  ls_residual?:     Scalar,        -- GMRES/FGMRES least-squares residual estimate
  breakdown_token?: BreakdownTag   -- a guarded-quantity validity tag
}
\end{lstlisting}

Every field is run-time and \emph{derived}: each is a pure function of the
post-step \code{Krylov} bundle---the residual norm read off the recurrence
scalar, the least-squares residual read off the rotated right-hand side, the
breakdown token read off a guarded dot product. And every optional field is
\emph{demand-prunable}: if no downstream consumer reads \code{ls\_residual}, the
kernel never computes it. The record exists precisely to make that pruning
structural---by typing the derived observations separately from the \code{Krylov}
state they are derived from, the calculus can state ``unread output, uncomputed''
as a law (\cref{ch:fold}) rather than a comment. \code{StepOutputs} is the
\emph{observation} a step makes about itself; do not confuse it with the
persistent \code{SimState.final\_res} statistic it eventually feeds.

\subsection{The five records, at a glance}

\Cref{tab:records} collects all five with their strata. Read it as the answer
to ``where does each piece of a solve's data live, and for how long.''

\begin{table}[t]
\centering
\caption[The five workhorse records and their strata]{The five records of a
Krylov-shaped solve, each with its stratum and lifetime. \code{OpParams} and
\code{SimState} are the two persistent strata (build-time vs.\ run-time);
\code{Krylov} and \code{PrevCarry} are the two short-lived ones; \code{SolveResult}
and its \code{StepOutputs} field are the result-side bundles a single step
returns.}
\label{tab:records}
\small
\setlength{\tabcolsep}{5pt}
\renewcommand{\arraystretch}{1.25}
\begin{tabular}{@{}p{24mm}p{20mm}p{34mm}p{30mm}@{}}
\toprule
Record & Stratum & Lifetime & Role \\
\midrule
\code{OpParams} & 1 (build) & whole solve, read-only & configuration + surfaces \\
\code{SimState} & 2 (run) & whole solve, evolving & caller-visible state \\
\code{Krylov} & 3 (ephemeral) & one restart cycle & algorithm workspace \\
\code{PrevCarry} & 4 (carry) & one inner sweep, reborn & threaded recurrence scalar \\
\code{StepOutputs} & run, derived & one step, demand-pruned & per-step readouts \\
\code{SolveResult} & result bundle & one step & what a step returns \\
\bottomrule
\end{tabular}
\end{table}

\section{The common stratification mistakes}
\label{sec:mistakes}

Each stratum boundary is a place where a value can be filed wrong, and each
mis-filing is a recognizable bug. We teach them as pitfalls because they are the
errors a reader coming from the mutable, instance-field world will reach for by
habit---the C++ stores everything as instance fields, so the original layout
offers no guidance on where a value \emph{belongs}.

\begin{pitfall}
\textbf{Ephemeral workspace in \code{SimState}.} The tempting mistake is to put
the basis \code{V} or the direction \code{p} into \code{SimState} because the
original C++ stored it as an instance field. But the instance-field layout is a
\emph{memory-reuse optimisation}, not a statement of lifetime: the field is
reused across restarts because allocating it once is cheap, not because it
\emph{lives} that long. Mathematically the basis dies at each restart. Put it in
\code{SimState} and you have claimed it survives the whole solve---so it
persists, wrongly, past the restart that should have discarded it, and a later
step can read a stale basis vector. The L4 type should reflect the mathematical
lifetime, not the allocation strategy. \emph{Workspace is stratum 3, never
stratum 2.}
\end{pitfall}

\Cref{fig:bug} draws this bug: a basis vector from before a restart, wrongly
surviving inside \code{SimState} and read by a step that should see only the
reborn bundle.

\begin{figure}[t]
\centering
\begin{tikzpicture}[font=\footnotesize, x=1mm, y=1mm]
  % restart boundary
  \draw[simRust, very thick, densely dashed] (50,-12) -- (50,20);
  \node[font=\scriptsize\bfseries, text=simRust] at (50,22) {restart};
  \node[font=\scriptsize, text=simGray] at (24,22) {cycle $r$};
  \node[font=\scriptsize, text=simGray] at (78,22) {cycle $r{+}1$};
  % SimState band wrongly carrying V across
  \fill[simBgBlue] (6,8) rectangle (104,16);
  \draw[simBlue,thick] (6,8) rectangle (104,16);
  \node[anchor=west,font=\scriptsize] at (8,12) {\code{SimState} \emph{(correct: x, it, ...)}};
  % the wrongly-persisted V, drawn straddling the restart line
  \fill[simBgRust] (30,-4) rectangle (70,4);
  \draw[simRust,thick] (30,-4) rectangle (70,4);
  \node[font=\scriptsize, text=simRust] at (50,0) {\code{V} wrongly kept in \code{SimState}};
  % the read that goes wrong
  \draw[backflow] (78,0) -- node[below,font=\tiny,text=simRust]{reads stale basis} (62,0);
  \node[opbox, minimum width=12mm] at (82,0) {step};
  % correct: Krylov reborn
  \node[font=\scriptsize, text=simGreen, align=center] at (24,-9)
    {\emph{correct:} \code{Krylov}\\dies at restart};
  \node[font=\scriptsize, text=simGreen, align=center] at (84,-9)
    {\emph{correct:} fresh\\\code{Krylov} born};
\end{tikzpicture}
\caption[The ephemeral-in-\code{SimState} bug]{Putting the ephemeral basis
\code{V} into \code{SimState} claims it survives the whole solve. It then
straddles the restart boundary (dashed), and a step in cycle $r{+}1$ reads a
basis vector from cycle $r$---stale data. The fix is to keep \code{V} in the
\code{Krylov} stratum, which is reborn at the restart, so there is no surviving
value to read across the boundary.}
\label{fig:bug}
\end{figure}

\begin{pitfall}
\textbf{Variant selectors in \code{SimState}.} The selectors \code{pc\_side},
\code{gs\_orthog}, \code{flexible} are \emph{read} during the solve, so it is
tempting to file them with the run-time state. They are build-time
configuration. They are read only by the constructed-operator surfaces---which
were built once, at construction, with the selectors already baked in---never by
the per-step kernel. File them in \code{SimState} and you have invited the kernel
to branch on them, which destroys variant absorption (\cref{sec:visible}): the
whole point was that the kernel \emph{cannot} see a raw selector. \emph{Selectors
are stratum 1, never stratum 2}---``read during the solve'' is not the same as
``run-time state.''
\end{pitfall}

\begin{pitfall}
\textbf{Forgetting that the carry is reborn at restart (or at each call).} The
recurrence carry \code{PrevCarry} threads a scalar across inner steps, so it
looks persistent. But it is reborn at each top-level call: the next solve
restarts its recurrence from the base case. Two failures follow from forgetting
this. \emph{If you file the carry in \code{OpParams}}, it persists between calls
and the next solve starts with last solve's $\beta_{k-1}$---a corrupted
recurrence. \emph{If your form reads a \code{Krylov} field across a restart
boundary} expecting the recurrence to continue, it reads from a bundle that was
discarded. The carry's lifetime is exactly ``one call's inner steps''---wider
than an ordinary temporary, narrower than configuration---and a form that gets
that window wrong is broken even when it type-checks against the field names.
\end{pitfall}

\begin{workedexample}{Build-time versus run-time in a preconditioned solve}{buildrun}
We take a preconditioned solve and split its work into the two great strata,
then show what each \emph{mislabeling} costs. This is the build-time/run-time
line of \cref{fig:buildrun} made concrete on one operation: the preconditioner.

\textbf{The two halves of a preconditioner.} A preconditioner $\mathsf{B}$
enters a solve in two distinct acts. First it is \emph{constructed}: for an
incomplete-factorization preconditioner this means \emph{factoring}---computing
the sparse triangular factors once, an expensive setup. Then, every iteration, it
is \emph{applied}: a triangular solve against the current residual, $\vz \gets
\mathsf{B}\vr$, cheap per call but run every step.

\textbf{The correct split.}
\begin{itemize}[nosep]
  \item \textbf{Build-time, once:} factor $\mathsf{B}$. The factors are
  \code{OpParams}-stratum data---configuration of the operator, captured at
  construction, read (applied) but never re-factored during the loop. In the
  monad they sit in the \code{setup} prologue.
  \item \textbf{Run-time, every step:} apply $\mathsf{B}$ to the residual. This
  is the per-iteration work; it touches \code{SimState} (the residual it reads,
  the iterate it helps advance) and lifts, in \cref{ch:tensors}, to a
  tensor-field expression.
\end{itemize}

\tcblower
\textbf{What each mislabeling costs.} The two halves are not interchangeable, and
swapping a label is a real failure---one wrong each way.

\begin{itemize}[nosep]
  \item \textbf{Calling the factorization run-time} (re-factoring $\mathsf{B}$
  every iteration) is \emph{wasted work}: you recompute, at full setup cost, a
  result that never changed. A solve that should factor once and apply a thousand
  times instead factors a thousand times. The answer is right; the cost is
  catastrophic. This is the classic ``loop-invariant computation left inside the
  loop.''
  \item \textbf{Calling the apply build-time} (applying $\mathsf{B}$ once, before
  the loop, and reusing the result) gives \emph{wrong results}: the apply depends
  on the current residual, which changes every step, so a single hoisted apply
  feeds every iteration a stale preconditioned vector and the solve does not
  converge to the right answer---if it converges at all.
\end{itemize}

\textbf{The lesson.} The build-time/run-time line is not cosmetic. Mislabel it
one way and you waste work; mislabel it the other and you compute the wrong
thing. The discriminating question is precisely \emph{lifetime}: does the value
change across iterations? The factorization does not (build-time); the apply's
result does (run-time). The matrix-free operator of \cref{ch:matrixfree} draws
exactly this line as its setup-versus-apply split, and for exactly this reason.
\end{workedexample}

\section{Lifetimes as the backbone of the design}

We close by collecting the thread. A solver's data is not one undifferentiated
``state.'' It is four kinds, sorted by lifetime: the build-time configuration
that is read for the whole solve, the run-time state that evolves across it, the
ephemeral workspace that lives one restart cycle, and the recurrence carry that
threads one inner sweep and is reborn each call. Each lifetime gets its own
record---\code{OpParams}, \code{SimState}, \code{Krylov}, \code{PrevCarry}---and
the result bundle \code{SolveResult} (with its demand-prunable \code{StepOutputs})
is what a single step returns.

The payoff is that the strata make the design's hardest questions structural
rather than disciplinary. You cannot persist ephemeral data past its window,
because the window's end is a fresh value, not an overwrite. You cannot let the
kernel branch on a variant, because the variant selectors are read-only and
reachable only through the constructed surfaces. And you cannot lose track of
which work runs once versus every step, because the build-time/run-time line is
the line between the \code{setup} prologue and the loop body---and, not
incidentally, the line between the work that prepares the loop and the work that
lifts to the tensor-field form of \cref{ch:tensors}. The variant-handling
machinery of \cref{ch:operators} rests entirely on this: with the strata in
place, absorbing a family of algorithms into one kernel is a typing invariant the
calculus checks, not a discipline the programmer must keep.

\summarysection
Every piece of a solver's state has a \emph{lifetime}, and classifying state by
lifetime---its \emph{stratum}---is what makes the design tractable. There are
four strata. \textbf{Build-time configuration} (\code{OpParams}) is captured once
at construction and read, read-only, for the whole solve; it holds the operator
surfaces, tolerances, and variant selectors, in a two-layer shape where the
kernel reaches the selectors only through constructed surfaces. \textbf{Run-time
state} (\code{SimState}) is the caller-visible bundle---\code{x}, \code{it},
\code{converged}, \code{final\_res}, \code{initial\_res}---that evolves every step
(as fresh values, never mutated) and is threaded by the \code{Solve} monad
(\cref{ch:monad}); its five-field shape is uniform across CG, GMRES, FGMRES, and
Chebyshev. \textbf{Ephemeral workspace} (\code{Krylov}: basis \code{V},
direction \code{p}) lives one restart cycle, is reborn at each restart, and is
threaded as a plain value, not in the monad---its lifetime defeats monadic
encoding. The \textbf{recurrence carry} (\code{PrevCarry}: \code{beta\_prev}) is
a fourth stratum, threaded across an inner loop but reborn each call. Making
lifetimes visible buys two guarantees: ephemeral data cannot be persisted past
its window, and variant absorption (\cref{ch:operators}) becomes a typing
invariant. The coarse \textbf{build-time versus run-time} split---setup once,
loop every step---is the same line the matrix-free operator draws
(\cref{ch:matrixfree}); only run-time work lifts to the tensor-field form
(\cref{ch:tensors}). The three classic mistakes are filing ephemeral workspace,
variant selectors, or the recurrence carry in the wrong stratum, and each is a
distinct, real bug.

\furtherreading
The conjugate-gradient and GMRES algorithms whose state we stratified, including
the restart structure that gives the ephemeral workspace its one-cycle lifetime
and the direction-update recurrence that gives the carry its cross-iteration
dependence, are developed in Saad's standard treatment of iterative
methods~\citep{saad2003}; reading the GMRES restart loop there with the four
strata in mind is the best exercise this chapter can recommend. The discipline of
making lifetime and effect visible in a value's type---so that ``what survives,
what is scratch, what threads'' is a fact the type states rather than a comment
the programmer leaves---is the contribution of pure functional programming and
its state monad; Peyton Jones's account of the lazy, purely functional
language~\citep{peytonjones2003} is where the threaded-state-as-a-value idea this
chapter applies to solvers is laid out in full, and it pairs naturally with the
monad of \cref{ch:monad}.

\exercisessection
\begin{enumerate}[nosep]
  \item For each of the four strata, state its lifetime in one phrase and name one
  field of the framework's record that lives there. (Example: ephemeral
  workspace---``one restart cycle''---\code{Krylov.V}.)
  \item A colleague proposes storing the iteration counter \code{it} in
  \code{OpParams} ``because it is just one integer and lives the whole solve.''
  Give the one reason this is wrong, and say which stratum \code{it} belongs to.
  (Hint: read-only versus evolved.)
  \item Take a GMRES solve and sort each of the following into its stratum: the
  preconditioner $\mathsf{B}$; the Krylov basis \code{V}; the Hessenberg matrix
  \code{H}; the iteration count \code{it}; the relative tolerance; the
  least-squares residual estimate $|s[j{+}1]|$. For the last one, say whether it
  is state or a derived \emph{output}, and which record it lives in.
  \item Explain, in two sentences, why the ephemeral \code{Krylov} bundle is
  threaded as a plain value rather than through the \code{Solve} monad. Your
  answer should turn on the word ``restart.''
  \item The chapter claims you ``cannot accidentally persist ephemeral data past
  its restart'' once it is in the \code{Krylov} stratum. Explain why the immutable
  fresh-value discipline of \cref{ch:records} is what makes this true, and what
  would go wrong in a mutable instance-field implementation.
  \item Why does putting a variant selector such as \code{pc\_side} into
  \code{SimState} defeat variant absorption? Trace the argument from ``the kernel
  can now read \code{pc\_side}'' to ``the kernel can now branch on the variant.''
  \item Distinguish the conjugate-gradient search direction $\vp$ (ephemeral,
  stratum 3) from the recurrence scalar $\beta_{k-1}$ (carry, stratum 4). Both
  ``come from the previous step''---state the precise question that separates
  them, and answer it for each.
  \item Consider a preconditioner that is \emph{rebuilt} partway through a solve
  (for instance, when the operator changes mid-sweep). Does the factorization
  stay build-time? Argue both sides, then propose how the stratification should
  describe a value that is mostly build-time but occasionally refreshed. (This is
  a genuine edge case; there is no single right answer, but a good answer names
  the lifetime precisely.)
\end{enumerate}
